% =====================================================================
%  04-arquitectura.tex - Estructura del workspace y andamiaje
% =====================================================================

\section{Estructura del workspace y andamiaje}\label{sec:impl-workspace}

{\emergencystretch=2em El workspace de Cargo contiene cinco crates del backend bajo \rutarepo{crates/}: \texttt{domain}, \texttt{application}, \texttt{infrastructure}, \texttt{web} y \texttt{bin}, conforme a la estructura de \cref{sec:workspace-cargo}. La dirección de dependencias apunta al dominio: \texttt{domain} no depende de otro crate del workspace; \texttt{application} depende de \texttt{domain}; \texttt{infrastructure} y \texttt{web} dependen de las dos capas interiores; y \texttt{bin}, raíz de composición que produce \texttt{despacho-cli} y \texttt{qadra-ca}, depende de todas. Este último es el único crate que utiliza \texttt{anyhow} en la frontera del proceso. HTTP consume los puertos de aplicación \texttt{Identity\allowbreak Workflow}, \texttt{Case\allowbreak Workflow}, \texttt{Case\allowbreak Document\allowbreak Workflow} y \texttt{Participant\allowbreak Workflow}, sin conocer los adaptadores concretos. El sexto miembro, \texttt{browser-crypto}, compila a WebAssembly y no depende de esos cinco crates. La \cref{tab:impl-crates} sintetiza las responsabilidades; las pruebas unitarias viven en \texttt{src} y las de integración en \texttt{tests/}. La convención de mantenimiento limita a menos de 400 líneas cada archivo de código, según \texttt{CONTRIBUTING.md}.\par}

\begin{table}[H]
  \centering
  \caption{Responsabilidades de los crates del workspace.}
  \label{tab:impl-crates}
  \begin{tabularx}{\linewidth}{@{}l X X@{}}
    \toprule
    \textbf{Crate} & \textbf{Responsabilidad} & \textbf{Elementos del flujo integrado} \\
    \midrule
    \texttt{domain} & Valores, reglas, puertos y errores; sin E/S & Identidad, expedientes, participantes, permisos, criptografía y auditoría. \\
    \texttt{application} & Casos de uso y fronteras de persistencia & Autenticación, asignaciones, directorio y documentos por expediente. \\
    \texttt{infrastructure} & Adaptadores de servicios y almacenamiento & PostgreSQL, Redis, cifrado, firma, TSA, ZIP e importación local. \\
    \texttt{web} & Adaptador HTTP y DTO & Extracción de peticiones, rutas, errores y límites de concurrencia. \\
    \texttt{bin} & Composición y comandos & Servidor, preparación de esquema, importación y herramientas offline. \\
    \texttt{browser-\allowbreak crypto} & Criptograf\'ia local del navegador & Claves cifradas, solicitudes CSR y firmas de bytes preparados. \\
    \bottomrule
  \end{tabularx}
\end{table}

La regla de dependencias no descansa en la disciplina del equipo sino en dos mecanismos verificables. El primero es el propio grafo de Cargo: como se argumentó en el diseño (\cref{sec:workspace-cargo}), una importación prohibida produce un error de compilación, no una advertencia de revisión. El segundo es una prueba de arquitectura, \rutarepo{crates/bin/tests/architecture.rs}, que lee el manifiesto del crate \texttt{domain} y falla si este declara dependencia alguna de \texttt{application}, \texttt{infrastructure} o \texttt{web}; esta prueba protege contra la degradación gradual del manifiesto, que el compilador no reportaría mientras los símbolos prohibidos no se usen. Una segunda prueba estructural, \rutarepo{crates/bin/tests/port_coverage.rs}, recorre el workspace y exige que cada \texttt{pub trait} del dominio ---cada puerto--- tenga al menos un adaptador (\texttt{impl <Trait> for}) en \texttt{infrastructure} y al menos un mock generado con \texttt{mockall} en los tests; para detectar un escáner degradado, la prueba usa puertos centinela (\texttt{AuditLog}, \texttt{Document\allowbreak Hasher}, \texttt{KeyManager}) cuya ausencia delataría un análisis roto.

La versión mínima soportada del compilador es Rust 1.88, declarada mediante \texttt{rust-version} en el \texttt{Cargo.toml} raíz y comprobada por un job obligatorio de integración continua. La elevación desde 1.78 permite adoptar correcciones de seguridad de \texttt{time} y de la familia PostgreSQL, conforme a \rutarepo{docs/adr/0013-raise-msrv-for-security-fixes.md}. Se conserva la familia de APIs construida alrededor de \texttt{der}~0.7 ---\texttt{x509-cert}~0.2, \texttt{cms}~0.2, \texttt{x509-tsp}~0.1 y \texttt{rsa}~0.9--- porque su sustitución exige una migración coordinada, no porque 1.78 siga siendo el requisito vigente. Las dependencias compartidas se declaran en \texttt{[workspace.dependencies]}. El perfil de release aplica \texttt{strip = true} y \texttt{lto = true}; en desarrollo se optimizan los núcleos \texttt{argon2}, \texttt{blake2}, \texttt{num-bigint-dig} y \texttt{rsa}, manteniendo información de depuración en el resto del workspace.

El archivo de dependencias bloqueadas fija \texttt{chacha20}~0.10.2 en sustitución de la versión 0.10.1 retirada del índice. RustCrypto documenta en esa versión la corrección del uso de una instrucción SSE4.1 en el backend SSE2 de las variantes de generador aleatorio y contador de 64 bits \cite{rustcrypto_chacha2026}. La actualización queda registrada en \rutarepo{Cargo.lock} y la comprobación de dependencias forma parte del corte reproducido del directorio.

La configuración se carga desde variables de entorno con \texttt{dotenvy} (archivo \texttt{.env} local, nunca versionado; la plantilla comentada \texttt{.env.example} sí lo está). La \cref{tab:impl-env} lista las variables. Dos decisiones de manejo de secretos son transversales: la carga de configuración registra en la bitácora únicamente la \emph{presencia} de cada variable, jamás su valor; y todo material secreto leído del entorno o de la entrada estándar (la KEK, la credencial del proveedor de sellado, las contraseñas) aterriza en buffers \texttt{Zeroizing} que sobrescriben su contenido al liberarse, conforme al ADR \rutarepo{docs/adr/0003-zeroize-secret-material-in-domain.md}.

\begin{table}[H]
  \centering
  \caption{Variables de entorno del prototipo (plantilla \texttt{.env.example}).}
  \label{tab:impl-env}
  \begin{tabularx}{\linewidth}{@{}l X@{}}
    \toprule
    \textbf{Variable} & \textbf{Propósito} \\
    \midrule
    \texttt{RUST\_LOG} & Filtro de severidad de la telemetría (por defecto \texttt{info}). \\
    \texttt{CINCEL\_BASE\_URL} & URL base del proveedor remoto de sellado de tiempo. \\
    \texttt{CINCEL\_API\_KEY} & Credencial del proveedor de sellado (secreto; se llena localmente). \\
    \texttt{DATABASE\_URL} & PostgreSQL para usuarios, expedientes, documentos y auditoría; conexión administrativa para migrar y rol restringido para servir. \\
    \texttt{REDIS\_URL} & Cadena de conexión Redis para sesiones, desafíos y límites. \\
    \texttt{KEK\_BASE64} & Clave de cifrado de claves: 32 bytes aleatorios en base64 (secreto). \\
    \texttt{NEW\_KEK\_BASE64} & KEK de reemplazo; la lee únicamente \texttt{vault rotate-kek}. \\
    \texttt{AUDIT\_LOG\_PATH} & Ruta de la bitácora offline; HTTP usa PostgreSQL (por defecto \texttt{audit-log.jsonl}). \\
    \texttt{PKI\_CA\_DIR} & Directorio de trabajo de la CA interna (por defecto \texttt{pki-ca}). \\
    \texttt{TSA\_DIR} & Directorio de trabajo de la TSA local (por defecto \texttt{pki-tsa}). \\
    \bottomrule
  \end{tabularx}
\end{table}

La telemetría usa \texttt{tracing} y \texttt{tracing-subscriber} con filtro por entorno, materializando la observabilidad estructurada prevista en el diseño (\cref{sec:observabilidad}). Todos los diagnósticos se emiten por la salida de error estándar, de modo que la salida estándar queda limpia para resultados y tuberías de shell. La promesa de que ningún secreto aparece en los diagnósticos no es una convención: la prueba \rutarepo{crates/bin/tests/log_redaction.rs} ejecuta los flujos con el nivel de log más verboso e inspecciona lo emitido, fallando si algún material secreto conocido aparece en la salida de error.

El workflow de integración continua ejecuta siete jobs en cada push y pull request (\cref{tab:impl-ci}). La comprobación de MSRV es obligatoria y usa Rust 1.88. Los jobs de pruebas y cobertura preparan PostgreSQL y Redis y separan las bases de identidad, expedientes y documentos, de modo que las pruebas de persistencia no se omiten por ausencia de configuración. El umbral de \rutarepo{scripts/coverage-gate.sh} exige 90\,\% de líneas por crate en \texttt{domain}, \texttt{application} e \texttt{infrastructure}; \texttt{bin} se reporta sin bloquear. Estos valores son criterios de aceptación configurados, no mediciones de una ejecución concreta.

\begin{table}[H]
  \centering
  \caption{Jobs del workflow de integración continua.}
  \label{tab:impl-ci}
  \begin{tabularx}{\linewidth}{@{}l l X@{}}
    \toprule
    \textbf{Job} & \textbf{Herramienta} & \textbf{Criterio de aprobación} \\
    \midrule
    format & \texttt{cargo fmt} & Formato canónico en todo el workspace. \\
    lint & \texttt{cargo clippy} & Cero advertencias (\texttt{-D warnings}). \\
    test & \texttt{cargo test} & Suite completa del workspace en verde. \\
    msrv & \texttt{cargo check} & Compila con Rust 1.88; comprobación obligatoria. \\
    coverage & \texttt{cargo-llvm-cov} & $\geq$ 90\,\% de líneas por crate criptográfico. \\
    deny & \texttt{cargo-deny} & Avisos, licencias, duplicados y fuentes en regla. \\
    binary-size & \texttt{cargo build -{}-release} & \texttt{despacho-cli} $\leq$ 25~MB. \\
    \bottomrule
  \end{tabularx}
\end{table}

El diseño comprometió cinco fitness functions como verificaciones automatizadas de los atributos de calidad (\cref{sec:fitness-functions}) \cite{ford2023}. La \cref{tab:impl-fitness} muestra dónde vive cada una en el repositorio. Cuatro de las cinco diseñadas se implementaron directamente; la función que exige un span de \texttt{tracing} por handler HTTP sigue pendiente aun cuando los handlers ya existen, y en su lugar se incorporó el umbral de cobertura por crate como función de aptitud adicional.

\begin{table}[H]
  \centering
  \caption{Fitness functions implementadas y su ubicación en el repositorio.}
  \label{tab:impl-fitness}
  \begin{tabularx}{\linewidth}{@{}l L{6.4cm} X@{}}
    \toprule
    \textbf{Id} & \textbf{Propiedad verificada} & \textbf{Verificación implementada} \\
    \midrule
    FF-1 & El dominio no depende de capas externas & Grafo de Cargo en cada build y \rutarepo{crates/bin/tests/architecture.rs}. \\
    FF-2 & Ningún log contiene material secreto & \rutarepo{crates/bin/tests/log_redaction.rs} al nivel de log más verboso. \\
    FF-3 & Todo puerto tiene adaptador y mock & \rutarepo{crates/bin/tests/port_coverage.rs} con puertos centinela. \\
    FF-4 & Cobertura $\geq$ 90\,\% en crates criptográficos & \rutarepo{scripts/coverage-gate.sh} sobre el reporte de \texttt{cargo-llvm-cov}. \\
    FF-5 & Binario release $\leq$ 25~MB & Job \texttt{binary-size} del workflow de integración continua. \\
    \bottomrule
  \end{tabularx}
\end{table}

La política de dependencias revisa avisos, licencias, duplicados y fuentes del grafo. Se configura en \texttt{deny.toml} para \texttt{cargo-deny} y en \texttt{audit.toml} para \texttt{cargo-audit}. Permanece declarada la excepción RUSTSEC-2023-0071 \cite{rustsec2023_0071}, relativa a temporización en operaciones de clave privada del crate \texttt{rsa}. El razonamiento original de \rutarepo{docs/adr/0002-rsa-signing-crate-and-advisory.md} suponía firma por CLI; esa premisa no describe la exposición HTTP actual. La autenticación y los límites de concurrencia acotan quién puede solicitar firmas y cuántas se ejecutan, pero no corrigen el canal lateral. Antes de un despliegue público se requiere revisar el modelo de amenaza y resolver la elección del firmante. La API local no se presenta como una configuración de producción validada.
